\chapter{Dasar Teori}

% ==========================================
% REGRESI
% ==========================================
\section{Regresi}
Regresi adalah metode statistik dalam \textit{machine learning} yang digunakan untuk memodelkan hubungan antara satu atau lebih variabel independen (fitur) dengan variabel dependen (target) yang bersifat numerik atau kontinu. Tujuan utama regresi adalah untuk memahami pola hubungan antara fitur dan target, memperkirakan nilai masa depan berdasarkan data sebelumnya, dan mengidentifikasi seberapa besar pengaruh masing-masing fitur terhadap output [\ref{modulPraktikum}, \ref{bishop}].

Dalam ekosistem \textit{machine learning}, regresi termasuk ke dalam kategori \textit{supervised learning}, karena model dilatih menggunakan pasangan \textit{input-output} sehingga bisa belajar memprediksi nilai target yang belum diketahui. Contoh penggunaan regresi antara lain yaitu prediksi harga rumah, estimasi permintaan barang, prediksi curah hujan, hingga evaluasi risiko. Regresi bekerja dengan cara menemukan fungsi atau model matematis yang paling sesuai dengan pola data. Model tersebut kemudian digunakan untuk memprediksi nilai baru berdasarkan *input* baru [\ref{geron}].

% ==========================================
% ALGORITMA REGRESI
% ==========================================
\section{Algoritma Regresi}
Algoritma regresi yang umum digunakan adalah \textit{Linear Regression}, \textit{Decision Tree Regression}, dan \textit{Random Forest Regression}. Masing-masing algoritma memiliki karakteristik dan cara kerja yang berbeda sesuai jenis data dan pola hubungan antar variabel.

\begin{enumerate}
    \item \textbf{Linear Regression}\\
    \textit{Linear Regression} adalah algoritma regresi paling dasar yang mengasumsikan bahwa hubungan antara fitur dan target bersifat linier. Model ini berusaha membangun sebuah garis atau \textit{hyperplane} terbaik yang meminimalkan \textit{error} antara prediksi dan nilai sebenarnya. Persamaan matematis dari \textit{Linear Regression} adalah sebagai berikut:
    \begin{equation}
        y = b_0 + b_1x_1 + b_2x_2 + \dots + b_nx_n
    \end{equation}
    Dimana:
    \begin{itemize}
        \item $y$ = target/prediksi
        \item $b_0$ = \textit{intercept} (nilai dasar)
        \item $b_n$ = koefisien regresi
        \item $x_n$ = nilai fitur
    \end{itemize}
    Model mencari koefisien terbaik menggunakan teknik \textit{Least Squares} (meminimalkan jumlah kuadrat *error*). Kelebihannya adalah sangat cepat, efisien untuk dataset besar, dan mudah diinterpretasikan. Namun kekurangannya adalah kurang cocok untuk pola non-linear dan sensitif terhadap \textit{outlier}.

    \item \textbf{Decision Tree Regression}\\
    \textit{Decision Tree Regression} bekerja dengan membangun struktur pohon keputusan sebagai representasi model. Algoritma ini membagi data menjadi beberapa subset berdasarkan titik *split* yang dianggap paling optimal. Setiap pembagian *node* dilakukan secara rekursif hingga mencapai \textit{leaf node} yang berisi nilai prediksi. Pemilihan *split* biasanya dilakukan berdasarkan pengukuran *error* seperti \textit{Mean Squared Error} (MSE) sehingga menghasilkan kelompok dengan variasi nilai target yang lebih homogen. Algoritma ini sangat efektif untuk hubungan non-linear dan tidak memerlukan normalisasi data, namun memiliki kelemahan utama berupa kecenderungan \textit{overfitting}.

    \item \textbf{Random Forest Regression}\\
    \textit{Random Forest Regression} merupakan pengembangan berbasis \textit{ensemble learning} yang membangun banyak pohon keputusan secara bersamaan. Setiap pohon dibangun menggunakan sampel data acak (\textit{bootstrap sampling}), serta subset fitur acak pada setiap pemisahan *node*. Hal ini menciptakan variasi pohon yang memberikan hasil prediksi stabil dan jauh lebih kebal terhadap *overfitting* dibanding pohon tunggal. Prediksi akhir diperoleh dari rata-rata hasil prediksi semua pohon. Selain itu, model ini dapat menampilkan \textit{feature importance} untuk mengetahui fitur yang paling berpengaruh.
\end{enumerate}

% ==========================================
% METRIK EVALUASI REGRESI
% ==========================================
\section{Metrik Evaluasi Regresi}
Evaluasi model regresi dilakukan menggunakan metrik khusus yang mampu mengukur kesalahan prediksi nilai numerik. Untuk menilai performa model regresi, digunakan dua metrik utama [\ref{modulPraktikum}]:

\subsection{Root Mean Squared Error (RMSE)}
RMSE mengukur rata-rata jarak antara nilai prediksi dan nilai sebenarnya dengan memberikan penalti lebih besar terhadap kesalahan yang besar. RMSE merupakan akar dari \textit{Mean Squared Error}, sehingga unitnya sama dengan unit target dan lebih mudah diinterpretasikan dalam masalah nyata. 
\begin{equation}
    RMSE = \sqrt{\frac{1}{n}\sum_{i=1}^{n}(y_{pred} - y_{true})^2}
\end{equation}
Nilai RMSE yang lebih kecil menandakan bahwa model menghasilkan prediksi yang lebih mendekati nilai sebenarnya (akurat). Karena RMSE memberikan penalti yang besar pada kesalahan ekstrem, metrik ini sensitif terhadap outlier, yang mana sangat cocok untuk mengevaluasi kasus regresi riil seperti fluktuasi harga rumah atau mobil.

\subsection{Coefficient of Determination ($R^2$)}
\textit{Coefficient of Determination} atau $R^2$ menunjukkan seberapa besar proporsi variansi target yang dapat dijelaskan oleh model. Nilai $R^2$ berada pada rentang matematis:
\begin{equation}
    0 \le R^2 \le 1
\end{equation}
Nilai $R^2$ yang semakin mendekati 1 menandakan bahwa model memiliki kemampuan prediksi yang sangat baik dan berhasil menangkap proporsi variansi pola hubungan data. Sebaliknya, nilai yang rendah berarti model kurang mampu memetakan hubungan antara input fitur dan target numerik.

\chapter{Hasil dan Pembahasan}

\section{Hasil Praktikum}

\subsection{Import Pustaka}
\begin{figure}[H]
    \centering
    \includegraphics[width=0.6\textwidth]{1.png}
    \caption{Source Code Import Pustaka}
    \label{fig:1}
\end{figure}
\begin{itemize}
    \item \textbf{Input:} Mengimpor pustaka dasar untuk pengolahan data dan visualisasi (\texttt{pandas}, \texttt{numpy}, \texttt{matplotlib.pyplot}, \texttt{seaborn}), fungsi preprocessing dari \texttt{sklearn} (\texttt{train\_test\_split}, \texttt{LabelEncoder}), algoritma pemodelan regresi (\texttt{LinearRegression}, \texttt{DecisionTreeRegressor}, \texttt{RandomForestRegressor}), serta fungsi metrik evaluasi (\texttt{mean\_squared\_error}, \texttt{r2\_score}, \texttt{mean\_absolute\_error}).
\end{itemize}

\subsection{Load Dataset}
\begin{figure}[H]
    \centering
    \includegraphics[width=0.7\textwidth]{2.png}
    \caption{Source Code Load Dataset}
    \label{fig:2}
\end{figure}
\begin{itemize}
    \item \textbf{Input:} Mengunduh dataset rumah (\texttt{kc\_house\_data.csv}) menggunakan \texttt{gdown}, lalu memuatnya ke dalam DataFrame \texttt{df} via \texttt{pd.read\_csv()}.
\end{itemize}

\subsection{Exploratory Data Analysis}
\begin{itemize}
    \item \textbf{EDA (Input \& Output):} 
    \begin{figure}[H]
        \centering
        \includegraphics[width=0.8\textwidth]{3.png}
        \caption{Output Exploratory Data Analysis}
        \label{fig:3}
    \end{figure}
    \begin{itemize}
        \item \textbf{df head:} Memanggil dan menampilkan 5 sampel baris teratas dari dataset harga rumah.
        \item \textbf{df tail:} Memanggil dan menampilkan 5 sampel baris terbawah.
        \item \textbf{df describe:} Memanggil dan mencetak ringkasan nilai statistik dari semua atribut numerik.
    \end{itemize}

    \item \textbf{Histogram (Input \& Output):}
    \begin{figure}[H]
        \centering
        \includegraphics[width=0.7\textwidth]{4.png}
        \caption{Plot Histogram}
        \label{fig:4}
    \end{figure}
    Menjalankan perintah pembuatan histogram dan menampilkan kumpulan grafik distribusi frekuensi dari tiap variabel numerik.

    \item \textbf{Matriks Korelasi (Input \& Output):}
    \begin{figure}[H]
        \centering
        \includegraphics[width=0.7\textwidth]{5.png}
        \caption{Matriks Korelasi}
        \label{fig:5}
    \end{figure}
    Menjalankan kode koreksi dan memunculkan tabel matriks korelasi untuk melihat fitur yang berpengaruh secara linear terhadap harga.

    \item \textbf{Pengecekan Missing Value (Input \& Output):}
    \begin{figure}[H]
        \centering
        \includegraphics[width=0.3\textwidth]{6.png}
        \caption{Pengecekan Missing Value}
        \label{fig:6}
    \end{figure}
    Menjalankan instruksi pencarian data kosong dan memunculkan deteksi 2 \textit{missing value} di atribut \texttt{sqft\_above}.

    \item \textbf{Pengecekan Atribut Kategorikal (Input \& Output):}
    \begin{figure}[H]
        \centering
        \includegraphics[width=0.6\textwidth]{7.png}
        \caption{Pengecekan Atribut Kategorikal}
        \label{fig:7}
    \end{figure}
    Menggunakan perintah \texttt{select\_dtypes} objek dan menampilkan daftar kosong, yang berarti tidak ada tipe data teks (semua murni numerik).
\end{itemize}

\subsection{Modelling}
\begin{itemize}
    \item \textbf{Modelling (Input \& Output):}
    \begin{figure}[H]
        \centering
        \includegraphics[width=0.8\textwidth]{8.png}
        \caption{Data Preprocessing}
        \label{fig:8}
    \end{figure}
    Melakukan imputasi median pada nilai kosong dan membagi dataset latih uji (70:30).
\end{itemize}

\subsubsection{Linear Regression}
\begin{itemize}
    \item \textbf{Linear Regression (Input \& Output):}
    \begin{figure}[H]
        \centering
        \includegraphics[width=0.7\textwidth]{9.png}
        \caption{Linear Regression Output}
        \label{fig:9}
    \end{figure}
    Mengaktifkan fungsi pelatihan \textit{Linear Regression} dan menampilkan nilai *r-squared* latih-uji, koefisien, dan *intercept*.

    \item \textbf{Evaluasi (Input \& Output):}
    \begin{figure}[H]
        \centering
        \includegraphics[width=0.5\textwidth]{10.png}
        \caption{Evaluasi Linear Regression}
        \label{fig:10}
    \end{figure}
    Melakukan kalkulasi kesalahan prediksi dan mencetak skor MSE, MAE, RMSE, serta perolehan $R^2$ sebesar 0.699.

    \item \textbf{Tes Visualisasi (Input):}
    \begin{figure}[H]
        \centering
        \includegraphics[width=0.6\textwidth]{11.png}
        \caption{Source Code Plotting LR}
        \label{fig:11}
    \end{figure}
    Mendeklarasikan penyusunan plot diagram bar untuk membandingkan deretan nilai aktual dengan prediksi hasil linear regresi.

    \item \textbf{Tes Visualisasi (Output):}
    \begin{figure}[H]
        \centering
        \includegraphics[width=0.9\textwidth]{12.png}
        \caption{Plot Linear Regression}
        \label{fig:12}
    \end{figure}
    Luaran grafik batang yang memvisualisasikan ketimpangan selisih tebakan regresi linear (oranye) terhadap data harga aslinya (biru).
\end{itemize}

\subsubsection{Decision Tree Regression}
\begin{itemize}
    \item \textbf{Decision Tree Regression (Input \& Output):}
    \begin{figure}[H]
        \centering
        \includegraphics[width=0.7\textwidth]{13.png}
        \caption{Decision Tree Regression}
        \label{fig:13}
    \end{figure}
    Mengeksekusi model \textit{Decision Tree} dengan batas \texttt{max\_depth=10}, dan mencetak angka performa akurasinya.

    \item \textbf{Evaluasi (Input \& Output):}
    \begin{figure}[H]
        \centering
        \includegraphics[width=0.5\textwidth]{14.png}
        \caption{Evaluasi Decision Tree}
        \label{fig:14}
    \end{figure}
    Memanggil metrik eror untuk mengevaluasi regresi Decision Tree dan mendapati kenaikan nilai $R^2$ pengujian menjadi 0.76.

    \item \textbf{Tes Visualisasi (Input):}
    \begin{figure}[H]
        \centering
        \includegraphics[width=0.6\textwidth]{15.png}
        \caption{Source Code Plotting DT}
        \label{fig:15}
    \end{figure}
    Source code persiapan rancangan diagram komparasi batang untuk model \textit{Decision Tree}.

    \item \textbf{Tes Visualisasi (Output):}
    \begin{figure}[H]
        \centering
        \includegraphics[width=0.9\textwidth]{16.png}
        \caption{Plot Decision Tree}
        \label{fig:16}
    \end{figure}
    Menampilkan plot tebakan (oranye) dan harga rill (biru) dari algoritma Decision Tree yang nampak sedikit membaik dibanding Linear.
\end{itemize}

\subsubsection{Random Forest Regression}
\begin{itemize}
    \item \textbf{Random Forest Regression (Input \& Output):}
    \begin{figure}[H]
        \centering
        \includegraphics[width=0.55\textwidth]{17.png}
        \caption{Random Forest Regression}
        \label{fig:17}
    \end{figure}
    Menjalankan perakitan algoritma \texttt{RandomForestRegressor} dan mencetak rasio latih-ujinya yang amat impresif (0.85).

    \item \textbf{Evaluasi (Input \& Output):}
    \begin{figure}[H]
        \centering
        \includegraphics[width=0.5\textwidth]{18.png}
        \caption{Evaluasi Random Forest}
        \label{fig:18}
    \end{figure}
    Perhitungan kerugian tingkat penyimpangan data yang menunjukkan *Random Forest* memiliki performa terbaik (RMSE terendah).

    \item \textbf{Tes Visualisasi (Input):}
    \begin{figure}[H]
        \centering
        \includegraphics[width=0.6\textwidth]{19.png}
        \caption{Source Code Plotting RF}
        \label{fig:19}
    \end{figure}
    Source code kerangka diagram untuk menge-plot data akhir dari regresi \textit{Random Forest}.

    \item \textbf{Tes Visualisasi (Output):}
    \begin{figure}[H]
        \centering
        \includegraphics[width=0.9\textwidth]{20.png}
        \caption{Plot Random Forest}
        \label{fig:20}
    \end{figure}
    Plot diagram komparasi bar yang menegaskan bahwa batang prediksi harga (oranye) sangat rapat dan presisi menutupi batang harga nyata (biru).
\end{itemize}


% ==========================================
% TUGAS PRAKTIKUM
% ==========================================
\clearpage
\section{Tugas Praktikum}

\subsection{Tujuan Penggunaan Dataset dan Definisi Atribut}
\textbf{Tujuan:} Dataset ini digunakan untuk melatih machine learning dalam memprediksi harga sebuah mobil. Dengan dataset ini, kita bisa membuat model regresi yang mampu memperkirakan nilai jual atau harga mobil di pasar Amerika Serikat berdasarkan spesifikasi bawaan mobil tersebut.

\textbf{Definisi Atribut:}
\begin{itemize}
    \item \textbf{Input (Variabel Independen):} Atribut yang mempengaruhi harga mobil, yang disimpan sebagai variabel X. Terdapat variabel bertipe numerik seperti panjang (\texttt{carlength}), lebar (\texttt{carwidth}), tinggi (\texttt{carheight}), bobot (\texttt{curbweight}), kapasitas mesin (\texttt{enginesize}), jumlah tenaga kuda (\texttt{horsepower}), serta konsumsi bahan bakar (\texttt{citympg} dan \texttt{highwaympg}). Ada juga variabel kategorikal (teks) seperti tipe bahan bakar (\texttt{fueltype}), tipe mesin, dan bentuk bodi.
    \item \textbf{Output (Variabel Dependen):} Atribut \texttt{price}, yaitu nilai jual dari mobil dalam mata uang USD. Nilainya yang berbentuk angka bebas dan berkelanjutan menandakan ini adalah penyelesaian kasus machine learning berjenis regresi.
\end{itemize}

\subsection{Variabel Paling Mempengaruhi Output (Korelasi Tinggi)}
Variabel yang paling mempengaruhi harga mobil (memiliki korelasi positif paling tinggi terhadap output \texttt{price}) adalah ukuran atau kapasitas mesin (\textbf{\texttt{enginesize}}). Semakin besar ukuran mesinnya, harga mobil semakin mahal. Di urutan selanjutnya, berat kendaraan (\texttt{curbweight}) dan daya kuda (\texttt{horsepower}) juga mempunyai nilai korelasi positif yang amat kuat.
Di sisi lain, atribut efisiensi bahan bakar di dalam kota (\texttt{citympg}) dan jalan raya (\texttt{highwaympg}) memiliki angka korelasi negatif yang paling tinggi, yang bermakna mobil yang per galonnya lebih irit akan cenderung lebih murah harganya.

\subsection{Implementasi Eksperimen Model Regresi}

\subsubsection{Import Pustaka}
\begin{figure}[H]
    \centering
    \includegraphics[width=0.8\textwidth]{T1.png}
    \caption{Import Pustaka}
    \label{fig:t1}
\end{figure}
\begin{itemize}
    \item \textbf{Input:} Memasukkan pustaka inti Python dan modul mesin \textit{sklearn}.
\end{itemize}

\subsubsection{Load Dataset}
\begin{figure}[H]
    \centering
    \includegraphics[width=0.8\textwidth]{T2.png}
    \caption{Load Dataset Car Price}
    \label{fig:t2}
\end{figure}
\begin{itemize}
    \item \textbf{Input \& Output:} Mengunduh set data \texttt{CarPrice\_Assignment} ke server dan menampilkannya dalam format *dataframe*.
\end{itemize}

\subsubsection{Exploratory Data Analysis}
\begin{itemize}
    \item \textbf{EDA (Input \& Output):} 
    \begin{figure}[H]
        \centering
        \includegraphics[width=0.9\textwidth]{T3.png}
        \caption{Pratinjau Data}
        \label{fig:t3}
    \end{figure}
    \begin{itemize}
        \item \textbf{df head:} Meninjau cuplikan data mobil pada 5 baris teratas.
        \item \textbf{df tail:} Meninjau cuplikan data pada 5 baris terbawah.
        \item \textbf{df describe:} Melihat tabel ringkasan standar deviasi dan kuartil.
    \end{itemize}

    \item \textbf{Histogram (Input \& Output):}
    \begin{figure}[H]
        \centering
        \includegraphics[width=0.6\textwidth]{T4.png}
        \caption{Histogram}
        \label{fig:t4}
    \end{figure}
    Mengonversi tiap data numerik dimensi mobil menjadi cetakan grafik distribusi histogram.

    \item \textbf{Matriks Korelasi (Input \& Output):}
    \begin{figure}[H]
        \centering
        \includegraphics[width=0.3\textwidth]{T5.png}
        \caption{Korelasi}
        \label{fig:t5}
    \end{figure}
    Eksekusi perhitungan matriks keterikatan korelasi tiap spek mesin mobil terhadap fluktuasi harganya.

    \item \textbf{Pengecekan Missing Value (Input \& Output):}
    \begin{figure}[H]
        \centering
        \includegraphics[width=0.7\textwidth]{T6.png}
        \caption{Cek Sel Kosong}
        \label{fig:t6}
    \end{figure}
    Mencari data bolong atau sel kosong, hasilnya memvalidasi bahwa set data ini telah utuh merata (angka 0).

    \item \textbf{Data Preprocessing (Input \& Output):}
    \begin{figure}[H]
        \centering
        \includegraphics[width=0.65\textwidth]{T7.png}
        \caption{Data Preprocessing}
        \label{fig:t7}
    \end{figure}
    \begin{itemize}
        \item \textbf{cek atribut kategorikal:} Pemisahan dan pelacakan daftar kolom bertipe string (teks).
        \item \textbf{label encoding:} Pengubahan deret kategori spesifikasi mobil yang masih berupa teks alfabet menjadi susunan representasi angka yang sah.
    \end{itemize}
\end{itemize}

\subsubsection{Modelling}
\begin{itemize}
    \item \textbf{Modelling (Input):}
    \begin{figure}[H]
        \centering
        \includegraphics[width=0.55\textwidth]{T8.png}
        \caption{Modelling (Setup)}
        \label{fig:t8}
    \end{figure}
    \begin{itemize}
        \item \textbf{def evaluate model:} Pembuatan cetak biru fungsi agar tidak perlu mengetikkan proses *fit*, prediksi, dan visualisasi berulang kali tiap mengetes jenis model regresi yang berbeda.
    \end{itemize}

    \item \textbf{Fungsi Regresi (Input):}
    \begin{figure}[H]
        \centering
        \includegraphics[width=0.8\textwidth]{T9.png}
        \caption{Deklarasi Fungsi}
        \label{fig:t9}
    \end{figure}
    \begin{itemize}
        \item \textbf{def run linear regression:} Kerangka pemanggilan dan inisiasi model *Linear Regression*.
        \item \textbf{def run decision tree:} Kerangka pemanggilan dan inisiasi model *Decision Tree*.
        \item \textbf{def run random forest:} Kerangka pemanggilan dan inisiasi model *Random Forest*.
    \end{itemize}

    \item \textbf{main (Input):}
    \begin{figure}[H]
        \centering
        \includegraphics[width=0.8\textwidth]{T10.png}
        \caption{Main Loop}
        \label{fig:t10}
    \end{figure}
    \begin{itemize}
        \item \textbf{memanggil fungsi:} Kode eksekutor penembakan 3 algoritma regresi secara simultan.
        \item \textbf{menampilkan table hasil evaluasi akhir:} Perintah transformasi hasil deret prediksi ke dalam kerangka *dataframe* pandas yang rapi.
    \end{itemize}

    \item \textbf{Hasil Menjalankan Linear Regression (Output):}
    \begin{figure}[H]
        \centering
        \includegraphics[width=0.9\textwidth]{T11.png}
        \caption{Eksekusi Linear}
        \label{fig:t11}
    \end{figure}
    \begin{itemize}
        \item \textbf{Hasil Menjalankan Linear Regression:} Informasi koefisien dari pola linear yang berhasil dilatih.
        \item \textbf{Evaluasi:} Catatan perbandingan rentetan 10 angka prediksi lawan nilai mobil aslinya.
        \item \textbf{Tes Visualisasi:} Diagram grafis perbandingan akurasi di mana batang prediksi (oranye) berselisih tajam dengan kenyataan (biru).
    \end{itemize}

    \item \textbf{Hasil Menjalankan Decision Tree Regression (Output):}
    \begin{figure}[H]
        \centering
        \includegraphics[width=0.9\textwidth]{T12.png}
        \caption{Eksekusi DT}
        \label{fig:t12}
    \end{figure}
    \begin{itemize}
        \item \textbf{Hasil Menjalankan Decision Tree Regression:} Bukti log terinisiasinya algoritma pohon keputusan.
        \item \textbf{Evaluasi:} Daftar kalkulasi perbandingan harga tebakan pohon dengan harga aktual pasar.
        \item \textbf{Tes Visualisasi:} Grafik batang *Decision Tree* yang pendaran oranyenya sedikit lebih menempel ketat dibanding regresi linear konvensional.
    \end{itemize}

    \item \textbf{Hasil Menjalankan Random Forest Regression (Output):}
    \begin{figure}[H]
        \centering
        \includegraphics[width=0.9\textwidth]{T13.png}
        \caption{Eksekusi RF}
        \label{fig:t13}
    \end{figure}
    \begin{itemize}
        \item \textbf{Hasil Menjalankan Random Forest Regression:} Bukti log eksekusi kumpulan *ensemble tree* (hutan acak).
        \item \textbf{Evaluasi:} Akurasi 10 tebakan awal model RF yang nyaris presisi mutlak tanpa melenceng jauh.
        \item \textbf{Tes Visualisasi:} Plot *bar chart* yang menyoroti bagaimana batang tebakan secara gemilang membayangi kurva batang aktualnya.
    \end{itemize}
\end{itemize}

\subsection{Tabel Performa dari Model Machine Learning}
\begin{figure}[H]
    \centering
    \includegraphics[width=0.9\textwidth]{T14.png}
    \caption{Table Comparison}
    \label{fig:t14}
\end{figure}
\begin{itemize}
    \item \textbf{Output:} Tabel final rekap pengujian yang merangkum persaingan komparasi metrik dari eksperimen:
    \begin{itemize}
        \item \textbf{Linear Regression, Decision Tree Regression, Random Forest Regression:} Deret nama algoritma pemodelan penantang.
        \item \textbf{$R^2$ Train, $R^2$ Test, MSE, RMSE, MAE:} Deret indikator rincian kalkulasi tingkat galat yang memvalidasi kemenangan telak *Random Forest*.
    \end{itemize}
\end{itemize}

\subsection{Model yang Paling Baik}
Berdasarkan bukti tabel metrik kinerja dari luaran gambar di atas, di antara ketiga algoritma regresi tersebut, \textbf{Random Forest Regression} merupakan model yang paling baik.

Alasan utamanya adalah:
\begin{itemize}
    \item \textbf{Tingkat Error Paling Kecil:} Dibandingkan Linear Regression dan Decision Tree, Random Forest memproduksi rentang nilai galat (yakni MSE, RMSE sebesar 2027.73, dan MAE) yang paling sedikit pada saat mengetes prediksi harga mobil, artinya tebakannya lebih tepat sasaran.
    \item \textbf{Skor ($R^2$) Paling Tinggi:} Model ini mencatatkan nilai \textit{R-Squared} yang jauh lebih tinggi di sisi data uji, yaitu 0.9407. Nilai yang mendekati 1 ini menjadi bukti nyata bahwa Random Forest sangat handal dan berhasil dalam menangkap pola harga mobil yang sebenarnya dari fitur-fiturnya.
    \item \textbf{Tahan Terhadap Overfitting:} Berbeda dengan kemunculan \textit{Decision Tree} tunggal yang menghasilkan $R^2$ latih nyaris sempurna (0.999) tetapi jeblok di pengujian (0.86), struktur model \textit{Random Forest} memadukan metode kumpulan pepohonan (\textit{ensemble}) yang menjaga keseimbangan model. Hal inilah yang mencegah algoritma tersebut menjadi \textit{overfit} pada dataset spesifikasi mobil yang begitu beraneka ragam parameternya.
\end{itemize}

\chapter{Kesimpulan}
\label{chap:kesimpulan}
Berdasarkan praktikum mengenai regresi pada \textit{machine learning} yang telah dilaksanakan—baik menggunakan dataset \textit{King County House Price} maupun dataset \textit{Car Price Prediction}—diperoleh beberapa simpulan penting sebagai berikut:

\begin{enumerate}
    \item \textbf{Pentingnya Data Preprocessing sebelum Pemodelan Regresi:} Tahapan \textit{Exploratory Data Analysis} (EDA) dan prapemrosesan (seperti imputasi nilai kosong/missing value, penanganan korelasi, serta \textit{label encoding} untuk fitur kategorial) terbukti sangat esensial. Pada kasus prediksi harga, memastikan data bersih dan terkonversi secara matematis menjadi kunci agar algoritma mesin bisa melakukan kalkulasi prediktif.
    
    \item \textbf{Performa Algoritma Regresi yang Beragam:} Dari komparasi ketiga model regresi yang diujikan (Linear Regression, Decision Tree Regression, dan Random Forest Regression), model \textit{ensemble} berupa \textbf{Random Forest Regression} terbukti menjadi model dengan performa yang paling unggul.
    
    \item \textbf{Analisis Metrik Evaluasi Regresi:} Keunggulan *Random Forest* tervalidasi secara faktual melalui kalkulasi metrik evaluasi. Pada pengujian prediksi harga mobil, algoritma ini mampu menekan kerugian galat ke tingkat paling rendah (RMSE 2027.73) dan memproduksi skor ketepatan regresi ($R^2$) tertinggi di atas algoritma pendahulu lainnya.
    
    \item \textbf{Resiliensi terhadap Overfitting:} Praktikum ini juga membuktikan kelemahan struktural model tunggal seperti *Decision Tree* yang sangat rentan terhadap *overfitting* (mencetak skor \textit{training} yang nyaris sempurna 0.99, namun jatuh di angka pengujian aktual). Pendekatan ansambel hutan acak (\textit{Random Forest}) sukses mengatasi rintangan varians tersebut dengan metode *bootstrap* dan *majority averaging*.
    
    \item \textbf{Korelasi Fitur Fisik terhadap Nilai Jual:} Dari segi keilmuan datanya, hasil eksplorasi \texttt{df.corr()} pada kedua studi kasus membuktikan bahwa spesifikasi dimensi mesin (seperti \texttt{enginesize} dan luas hunian \texttt{sqft\_living}) memiliki korelasi regresi berbanding lurus yang amat dominan sebagai faktor pendorong lonjakan harga sebuah aset.
\end{enumerate}
\chapter{Daftar Pustaka}
\begin{enumerate}
    \item \label{modulPraktikum} Dr. Eng. Ir. Ganjar Alfian, S.T., M. Eng., dan Tim Asisten. \textit{Modul Praktikum Penambangan Data: Machine Learning Regresi}. Yogyakarta: Departemen Teknik Elektro dan Informatika, Sekolah Vokasi, Universitas Gadjah Mada, 2026.
    
    \item \label{mckinney} W. McKinney. \textit{Python for Data Analysis: Data Wrangling with Pandas, NumPy, and Jupyter}, 3rd ed. Sebastopol, CA: O'Reilly Media, 2022.
    
    \item \label{geron} A. Géron. \textit{Hands-On Machine Learning with Scikit-Learn, Keras, and TensorFlow}, 3rd ed. Sebastopol, CA: O'Reilly Media, 2022.
    
    \item \label{bishop} C. M. Bishop. \textit{Pattern Recognition and Machine Learning}. New York, NY: Springer, 2006.
    
    \item \label{kaggleHouse} Shiva Chandel. \textit{KC House Data}. Kaggle, 2021. [Online]. Tersedia: \url{https://www.kaggle.com/datasets/shivachandel/kc-house-data}.
    
    \item \label{kaggleCar} Hellbuoy. \textit{Car Price Prediction}. Kaggle, 2020. [Online]. Tersedia: \url{https://www.kaggle.com/datasets/hellbuoy/car-price-prediction/data}.
    
    \item \label{colabLaporan} H. C. A. Paliling. \textit{S5-PPD-542159 (Source Code Laporan Praktikum)}. Google Colaboratory, 2026. [Online]. Tersedia: \url{https://colab.research.google.com/drive/1bvBJJTwzyzKHsH95cHhvi076hqc8DzNy?usp=sharing}.
\end{enumerate}
